\documentclass[11pt]{report}
\input{../pri-end/T0_preamble_standalone_A4_De}
% T0_preamble_patches.tex — LOKALER PLATZHALTER
% Im Repo durch die offizielle Version ersetzen.


\title{Dok.~312 -- Die Abschlussskala:\\
$\Lambda$ als abgeleitete Größe der kompakten Geometrie}
\author{Johann Pascher}
\date{August 2026}

\begin{document}
\maketitle
\tableofcontents

% ============================================================
\chapter*{Dok.~312 -- Die Abschlussskala}
\addcontentsline{toc}{chapter}{Dok.~312 -- Die Abschlussskala}
% ============================================================

% ============================================================
\section*{Worum es geht}
% ============================================================

Die kosmologische Konstante $\Lambda$ hat in den Einsteinschen
Feldgleichungen eine Funktion: Sie setzt die großskalige Krümmungsskala.
Diese Funktion braucht auch die FFGFT -- aber nicht als freie Konstante
der Gleichungen, sondern als \textbf{abgeleitete Größe}. Dieses Dokument
arbeitet die Identifikation aus:

\[
  \Lambda \;\longrightarrow\; \Lambda^{*} := \frac{1}{R_H^2}
  \qquad\text{mit}\qquad
  R_H = \frac{c}{H_0} = \frac{2}{\pi}\,\frac{\bar\lambda_e}{\xi^{10}}.
\]

Drei Ergebnisse:
\begin{enumerate}
  \item Die dimensionslose Form
        $\Lambda^{*}\bar\lambda_e^{2} = (\pi/2)^{2}\,\xi^{20}$
        ist eine reine Zahl -- der SI-Bezug steckt vollständig im
        Kammerton $\bar\lambda_e$ (Dok.~310).
  \item Das $10^{123}$-Feinabstimmungsproblem zerfällt in eine
        \emph{Strukturzerlegung}: $\sim\!122$~Dekaden
        $= 77{,}5$ (aus $\xi^{20}$) $+\,44{,}8$ (aus
        $(l_P/\bar\lambda_e)^2$) $-\,0{,}4$ (Ordnung-1). Aus einer
        Feinabstimmungsfrage wird die Strukturfrage
        \enquote{warum Stufe~10 der $\xi$-Leiter?}
  \item Die Vererbung aus Dok.~309 gilt unverändert: Solange der
        Exponent~10 nicht vorwärts hergeleitet ist (P20), ist
        $\Lambda^{*}$ [K]~\emph{bedingt} -- korrekt per Konstruktion,
        nicht per Ableitung.
\end{enumerate}

Das Dokument ist kein Teil der A-Serie. Es überschreibt keine
bestehenden Dokumente; notwendige Folgeänderungen werden separat
in Dok.~190 registriert.
Bezüge: Dok.~182 ($R_H$ als geometrische Maximalskala),
Dok.~190 (P20, P33), Dok.~267 (Entartung), Dok.~279 ($H_0$-Kette),
Dok.~308 ($a_0$, Trägheitsregime; A272-Vorläufer $\Lambda$-Lesart),
Dok.~309 (Vererbungsproblem), Dok.~310 (Kammerton-Prinzip).

% ============================================================
\section*{Leseführung: der rote Faden}
% ============================================================

Das Dokument ist gewachsen; zehn Kapitel und ein Anhang. Damit der
Bogen sichtbar bleibt, hier die Erzählung in Kurzform.

\textbf{Ausgangspunkt (Kap.~A--C): $\Lambda$ bekommt einen
Träger.} Die Feldgleichungen enthalten einen Term, der die größte
Krümmungsskala setzt; $\Lambda$CDM \emph{fittet} seinen Wert. Hier
wird er \emph{identifiziert}: Die FFGFT besitzt bereits eine größte
Länge, $R_H$ -- also $\Lambda^{*} = 1/R_H^2$, die Abschlussskala,
kein neuer Parameter. Daraus zerfällt das $10^{123}$-Problem in
Arithmetik zweier Strukturverhältnisse. Ehrlich bleibt: Alles hängt
am Exponenten~10 (P20), daher überall das bedingte Label [K$|$P20].

\textbf{Die Weichenstellung (Kap.~A.1): volle Übernahme der
Feldgleichungen.} Lovelock erzwingt ihre Form und lässt nur
$G$-Wert, $\Lambda$-Wert und Lösungswahl offen. Die FFGFT ersetzt
die Gleichungen also nicht -- sie liefert, was offen ist.

\textbf{Der Preis (Kap.~D, Anhang): die Stabilitätsfrage.} Voll
gültige Gleichungen heißen dynamische Geometrie -- warum rollt der
Torus weder ein noch aus (Eddington 1930)? Der Anhang rechnet den
Kandidaten: Windungs- gegen Impulsmoden stabilisieren genau bei
$R_H$, und die nötige Spannung ist wieder $\Lambda^{*}$ -- eine
Skala, zwei Rollen. Offen als Pflicht~C.

\textbf{Der Schwenk zur Zeit (drei Kapitel, eine Botschaft).}
\emph{Zeit in den Feldgleichungen:} Zeit ist Metrikrichtung, keine
äußere Uhr; die FFGFT sagt, woran der Takt hängt -- an der Masse.
\emph{Statik der Relation:} Zeitdilatation ist eine statische
Landkarte $m(x)$; wer sie dynamisch erzählt, vergisst die
Massenseite -- und Expansion vs.\ Massenlauf ist Rahmenwahl
(Wetterich), keine Messung. \emph{Denken in Lichtwegen:} Warum das
übersehen wird -- der Bote unserer Kosmologie trägt selbst keine
Uhr ($d\tau = 0$); Uhren sind Massen, und die entscheidenden Tests
sind Materie-Ablesungen. Dazu die Grenze bei $c$: relational, nicht
innerlich.

\textbf{Die Zuspitzung (kompakte Zeitrichtung).} Ist Zeit
Metrikrichtung und lebt alles auf T$^4$, ist auch die Zeitrichtung
kompakt. Ertrag: Zeitzyklus $92$~Gyr, Quant exakt $\hbar H_0$,
Temperatur exakt Gibbons--Hawking ohne Horizont, und daraus
$a = cH_0$ -- die $a_0$-Verankerung aus einem Argument statt
zweien. Preis: die Kausalitätsfrage, Pflicht~D.

\textbf{Die Erdung (Relativbewegung).} Wir kreisen trotzdem:
lokal Bewegung \emph{auf} der Landkarte (reine Kinematik), global
zeichnet die Topologie ein Ruhesystem aus -- und beobachtet wird
genau eines, der CMB-Dipol ($369{,}8$~km/s exakt). Konvention wird
Strukturaussage, prüfbar ($4{,}2'$-Versatz).

\textbf{Der Epilog (Einstein).} Er hatte alle Zutaten für
$\tilde T \cdot m = 1$, aber die Relation hätte $\hbar$ zum
Fundament gemacht -- den einen Stein, den er nicht verbauen wollte.
Die FFGFT verbaut ihn und übernimmt dafür seine Gleichungen
unverändert: Arbeitsteilung statt Widerlegung. Damit schließt der
Epilog an die Weichenstellung von Kap.~A an.

\medskip
\noindent\textbf{In einem Satz:} Einsteins Gleichungen gelten
ganz; was $\Lambda$CDM fittet ($\Lambda$) und was Einstein offen
ließ (woran die Zeit hängt), liefert die kompakte T$^4$-Struktur
aus einer einzigen Leiterstufe -- zum deklarierten Preis von vier
offenen Pflichten: A~(Exponent~10), B~($\kappa$), C~(Stabilität),
D~(Kausalität der kompakten Zeit).

% ============================================================
\section*{Konventionen}
% ============================================================

\noindent\textbf{Zeichenerklärung.}

\medskip
\begin{tabular}{@{}lp{11cm}}
\textbf{[K]}       & \emph{Kernableitung} -- mathematisch oder numerisch
                     aus den Grundsetzungen hergeleitet; reproduzierbar. \\[4pt]
\textbf{[K|P20]} & Kernableitung \emph{bedingt auf} die Setzung
                     des Exponenten~10 (P20). \\[4pt]
\textbf{[S]}       & \emph{Strukturaussage} -- begründet, aber noch nicht
                     vollständig bewiesen. \\[4pt]
\textbf{[SETZUNG]} & \emph{Setzung} -- deklarierte Annahme, nicht weiter
                     hergeleitet. \\
\end{tabular}

\medskip
\noindent\textbf{Zentrale Größen.}

\medskip
\begin{tabular}{@{}lp{10.5cm}}
$\xi = 4/30000$ & Einziger Fundamentalparameter [SETZUNG, P33]. \\[4pt]
$\bar\lambda_e$ & Reduzierte Compton-Wellenlänge des Elektrons,
                  $3{,}8616\times10^{-13}$~m; Kammerton (Dok.~310). \\[4pt]
$l_P$           & Planck-Länge, $1{,}6163\times10^{-35}$~m. \\[4pt]
$H_0$           & $(\pi/2)\,c\,\xi^{10}/\bar\lambda_e
                   = 66{,}82$~km/s/Mpc [K|P20]. \\[4pt]
$R_H = c/H_0$   & Geometrische Maximalskala der kompakten Geometrie,
                  $1{,}384\times10^{26}$~m -- kein Expansionshorizont
                  (Dok.~182). \\[4pt]
$\Lambda^{*}$   & Abschlussskala, $:= 1/R_H^2$; Träger der
                  $\Lambda$-Funktion in der FFGFT. \\[4pt]
$\kappa$        & Ordnung-1-Faktor der Beobachtungsrelation
                  $\Lambda_\text{obs} = \kappa\,\Lambda^{*}$; offen. \\
\end{tabular}

% ============================================================
\chapter{Die Funktion von $\Lambda$ -- und wer sie trägt}
\label{ch:funktion}
% ============================================================

\section{$\Lambda$ steht links}
\label{sec:links}

In $G_{\mu\nu} + \Lambda g_{\mu\nu} = (8\pi G/c^4)\,T_{\mu\nu}$ ist
$\Lambda$ ein Geometrieterm mit Dimension $1/\text{Länge}^2$. Seine
Funktion: eine großskalige Krümmungsskala setzen, gegen die alle
lokalen Krümmungen klein sind. Diese Funktion ist lesartunabhängig --
sie existiert in der Expansionslesart ebenso wie in der statischen.

\textbf{Die FFGFT übernimmt die Feldgleichungen als gültig.} Nach dem
Lovelock-Theorem ist $G_{\mu\nu} + \Lambda g_{\mu\nu}$ in $D=4$ der
einzige divergenzfreie Tensor zweiter Ordnung aus der Metrik -- die
\emph{Form} der Gleichungen ist erzwungen, der $\Lambda$-Term
zwingend zugelassen. Was die Gleichungen offen lassen, ist genau
dreierlei: der Wert von $G$, der Wert von $\Lambda$, und die Wahl
der Lösung, die auf den Kosmos angewendet wird. Modellabhängig ist
also nicht die Gleichung, sondern die $\Lambda$CDM-Trias
\{$\Lambda$-Fit, FLRW-Anwendung, Expansionsdeutung von $z$\}.
Das Programm der FFGFT ist entsprechend nicht, die Gleichungen neu
herzuleiten, sondern zu liefern, was sie offen lassen: den
$\Lambda$-Wert (dieses Dokument), den $G$-Wert (A145) und die
Begründung der Lösungswahl. Die Zwei-Hälften-Synthese
(Dok.~307/308) zeigt zusätzlich, dass die FFGFT den lokalen Sektor
\emph{unabhängig} reproduziert -- Konsistenz, nicht bloße
Kompatibilität.

\section{Die statische Lesart braucht keinen neuen Freiheitsgrad}

In der FFGFT existiert bereits eine großskalige Längenskala:
$R_H$, die geometrische Maximalskala der kompakten
T$^4$/Z$_3$-Geometrie (Dok.~182). Die Identifikation

\begin{equation}
  \Lambda^{*} := \frac{1}{R_H^2}
  \label{eq:def}
\end{equation}

führt \emph{keinen} neuen Parameter ein -- sie benennt eine schon
vorhandene Größe um ihrer Funktion nach. [K|P20]

\textbf{Abgrenzung zur A-Serie (A272):} Dort wurde gezeigt, dass
$\Lambda$ \emph{als fundamentale Naturkonstante} ein Artefakt der
Expansionslesart ist. Das bleibt bestehen. Dieses Dokument liefert die
konstruktive Ergänzung: Die \emph{Funktion} von $\Lambda$ ist real und
wird von einer abgeleiteten Skala getragen. Artefakt als Konstante --
real als abgeleitete Skala.

\section{Namensgebung}

Die Größe $\Lambda^{*}$ heißt in der FFGFT nicht
\enquote{kosmologische Konstante}, sondern \textbf{Abschlussskala}
(der kompakten Geometrie). Der Name folgt derselben Disziplin wie
$a_0$ als \enquote{Übergangsskala des Trägheitsregimes} (Dok.~308):
Er sagt, \emph{was} die Größe ist -- eine Stufe der $\xi$-Leiter,
keine neue Naturkonstante.

% ============================================================
\chapter{Die Kette [K|P20]}
% ============================================================

\section{Dimensionsbehaftete Form}

Aus $H_0 = (\pi/2)\,c\,\xi^{10}/\bar\lambda_e$ folgt unmittelbar:

\begin{equation}
  \Lambda^{*}
  = \frac{H_0^2}{c^2}
  = \left(\frac{\pi}{2}\right)^{2}\frac{\xi^{20}}{\bar\lambda_e^{2}}
  = 5{,}218\times10^{-53}~\text{m}^{-2}.
  \label{eq:kette}
\end{equation}

Vergleich mit der Beobachtung (Planck~2018,
$\Lambda_\text{obs} \approx 1{,}106\times10^{-52}$~m$^{-2}$):

\begin{equation}
  \kappa := \Lambda_\text{obs}\,R_H^{2} = 2{,}12.
  \label{eq:kappa}
\end{equation}

In $\Lambda$CDM ist dieser Faktor $3\,\Omega_\Lambda \approx 2{,}07$
(mit $H_0^\text{Planck}$) -- dort ein Fit. In der FFGFT ist $\kappa$
\textbf{offen und muss hergeleitet werden} (Kap.~\ref{ch:pflichten}). Er wird hier
ausdrücklich \emph{nicht} durch eine passende Zahlenkombination
geschlossen.

\section{Dimensionslose Form}

Multiplikation von Gl.~\eqref{eq:kette} mit $\bar\lambda_e^2$:

\begin{equation}
  \boxed{\;\Lambda^{*}\bar\lambda_e^{2}
  = \left(\frac{\pi}{2}\right)^{2}\xi^{20}
  = 7{,}78\times10^{-78}\;}
  \label{eq:dimlos}
\end{equation}

Eine reine Zahl. Der gesamte SI-Bezug steckt -- wie bei $H_0/\nu_e$
in Dok.~309 -- in der Wahl des Kammertons $\bar\lambda_e$. Die
Abschlussskala ist damit vollständig in die Resonanzgeometrie
(Dok.~310) eingeordnet: dimensionslose Struktur ($\xi^{20}$, $\pi/2$),
ein einziger Einheitenanker.

% ============================================================
\chapter{Auflösung der $10^{123}$-Feinabstimmung}
% ============================================================

\section{Die Standardformulierung}

Das Feinabstimmungsproblem wird üblicherweise als Verhältnis
$\rho_\Lambda/\rho_\text{QFT} \sim 10^{-123}$ formuliert, äquivalent

\[
  \Lambda_\text{obs}\,l_P^{2} \sim 10^{-122}.
\]

$\Lambda$CDM hat dafür keine Erklärung: Die Zahl wird gesetzt.

\section{Die Strukturzerlegung [K|P20]}

In der FFGFT ist dieselbe Zahl ein Produkt zweier Strukturverhältnisse:

\begin{equation}
  \Lambda^{*}\,l_P^{2}
  = \underbrace{\left(\frac{\pi}{2}\right)^{2}}_{-0{,}4~\text{dex}}
    \;\underbrace{\xi^{20}}_{-77{,}5~\text{dex}}
    \;\underbrace{\left(\frac{l_P}{\bar\lambda_e}\right)^{2}}_{-44{,}8~\text{dex}}
  = 1{,}36\times10^{-122}.
  \label{eq:zerlegung}
\end{equation}

\begin{center}
\begin{tabular}{lll}
\toprule
Faktor & Dekaden & Herkunft \\
\midrule
$\xi^{20}$ & $-77{,}5$ & Stufe~10 der $\xi$-Leiter, quadriert \\
$(l_P/\bar\lambda_e)^2$ & $-44{,}8$ & Kammerton-zu-Planck-Verhältnis \\
$(\pi/2)^2$ & $-0{,}4$\;($+0{,}39$) & Geometriefaktor der $H_0$-Kette \\
\midrule
Summe & $-121{,}9$ & $\Lambda^{*}l_P^2 = 10^{-121{,}87}$ \\
\bottomrule
\end{tabular}
\end{center}

\textbf{Interpretation:} Es gibt nichts abzustimmen. Die
$\sim\!122$~Dekaden sind kein unerklärter Zufall, sondern die
arithmetische Folge davon, \emph{auf welcher Stufe der Leiter} die
Abschlussskala sitzt und \emph{wo der Kammerton relativ zur
Planck-Skala} liegt. Die Feinabstimmungsfrage
(\enquote{warum ist $\Lambda$ so klein?}) wird zur Strukturfrage
(\enquote{warum Exponent~10?}) -- und genau diese Frage ist bereits
als P20 registriert. Die Zerlegung erzeugt also keine neue offene
Stelle; sie führt zwei scheinbar getrennte Probleme
($10^{123}$-Problem, P20) auf \emph{eine} zurück. [S]

\section{Was die Zerlegung nicht leistet}

Sie erklärt nicht, warum $n=10$. Solange P20 offen ist, gilt: Die
Auflösung der Feinabstimmung ist \emph{bedingt} -- sie steht und
fällt mit der Vorwärtsherleitung des Exponenten. Das ist dieselbe
Vererbung wie in Dok.~309 und wird nicht wegdiskutiert.

% ============================================================
\chapter{Einstein 1917 und die Stabilitätsfrage}
\label{ch:einstein1917}
% ============================================================

\section{Der historische Befund}

Einsteins statisches Universum (1917) verlangt exakt
$\Lambda = 1/R^2$ (mit $R$ als Krümmungsradius der geschlossenen
statischen Lösung). Dass die FFGFT-Identifikation
$\Lambda^{*} = 1/R_H^2$ formal genau dort landet, ist bemerkenswert --
aber die klassische statische Lösung ist \textbf{instabil}
(Eddington 1930): Eine kleine Störung lässt sie kollabieren oder
expandieren. Das war das historische Todesurteil der statischen
Lesart innerhalb der ART.

\section{Die Stabilitätsfrage ist offen -- und wird nicht umgangen}
\label{sec:stabfrage}

Da die FFGFT die Feldgleichungen als voll gültig übernimmt
(Abschn.~\ref{sec:links}), ist die Metrik dynamisch -- und damit greift die
Eddington-Frage grundsätzlich auch hier. Ein Ausweichen über
\enquote{die Geometrie ist nicht dynamisch} stünde im Widerspruch
zur eigenen Position und wird ausdrücklich \emph{nicht} in Anspruch
genommen.

Der relevante Unterschied zur klassischen Einstein-Lösung von 1917
liegt woanders: Einsteins statisches Universum ist eine Lösung mit
offener räumlicher 3-Sphäre-Geometrie, deren Gleichgewicht ein
instabiler Balancepunkt zwischen Anziehung und $\Lambda$-Term ist.
Die FFGFT-Konfiguration ist dagegen auf der kompakten
T$^4$/Z$_3$-Topologie definiert. \textbf{Kandidat für die
Stabilisierung ist die Topologie selbst}: kompakte Zyklen als
Randbedingung, die den Störungsmoden, welche die klassische Lösung
kollabieren oder expandieren lassen, nicht offenstehen -- analog
dazu, wie Windungszahlen auf kompakten Zyklen nicht stetig
deformierbar sind.

Status ehrlich: Das ist ein \emph{Kandidatenmechanismus} [S], kein
Nachweis. Der Vorwärtsbeweis müsste die Störungsanalyse der
statischen Konfiguration auf T$^4$/Z$_3$ durchführen und zeigen,
dass die instabilen Moden der Einstein-Lösung durch die kompakte
Topologie ausgeschlossen oder stabilisiert werden. Bis dahin steht
die Eddington-Frage als offene Stelle im Register (Kap.~\ref{ch:pflichten}, Punkt~C)
-- schärfer gestellt als zuvor, weil die volle Gültigkeit der
Feldgleichungen sie unausweichlich macht.

% ============================================================
\chapter{Querverankerung und lokaler Sektor}
% ============================================================

\section{Dieselbe Leiterstufe trägt $a_0$ [K|P20]}

\[
  c\,H_0 = 6{,}49\times10^{-10}~\text{m/s}^2,
  \qquad
  \frac{c\,H_0}{2\pi} = 1{,}03\times10^{-10}~\text{m/s}^2
  \quad\text{vs.}\quad
  a_0^\text{RAR} \approx 1{,}2\times10^{-10}~\text{m/s}^2.
\]

Abschlussskala ($\Lambda^{*}$) und Übergangsskala des
Trägheitsregimes ($a_0$, Dok.~308) sitzen auf \emph{derselben} Stufe
$\xi^{10}$. Ein reiner $H_0$-Fit würde das nicht erzwingen -- das ist
die Querverankerung aus Dok.~309, um einen Eintrag erweitert:

\begin{center}
\begin{tabular}{llll}
\toprule
Sektor & Größe & $\xi$-Abhängigkeit & Status \\
\midrule
Kosmologie & $H_0$ & $\xi^{10}$ & [K|P20] \\
Kosmologie & $\Lambda^{*} = 1/R_H^2$ & $\xi^{20}$ & [K|P20] \\
Trägheitsregime & $a_0$ & $\xi^{10}$ & [K]/[S] \\
Teilchenphysik & Leptonmassenleiter & $\xi$-Potenzen & [K] \\
Quantenmechanik & $K_\text{frak} = 1-100\xi$ & $\xi$ & [K] \\
\bottomrule
\end{tabular}
\end{center}

\section{Der lokale Sektor bleibt unberührt [K]}

$\Lambda^{*} \sim 10^{-52}$~m$^{-2}$ ist gegen jede lokale Krümmung
vernachlässigbar (Sonnensystem: Krümmungsskalen
$\gtrsim 10^{-27}$~m$^{-2}$ am Sonnenrand; Verhältnis
$\gtrsim 10^{25}$). Die Schwarzschild-Lösung mit $\Lambda = 0$, die
Zwei-Hälften-Synthese und alle klassischen Tests (Dok.~307/308)
bleiben exakt gültig -- die Abschlussskala liefert automatisch die
richtige Hierarchie, ohne dass sie lokal abgeschaltet werden müsste.

% ============================================================
\chapter{Was herzuleiten bleibt}
\label{ch:pflichten}
% ============================================================

\section{A -- Der Exponent 10 (P20, prioritär)}

Die Vorwärtsherleitung von $n=10$ aus der T$^4$/Z$_3$-Struktur ist
die eine Stelle, an der die gesamte Kette hängt. Ohne sie: alles
[K|P20], Vererbung nach Dok.~309.

\textbf{Kandidaten-Lesarten -- ausdrücklich als solche markiert:}
\begin{itemize}
  \item $10$ = Zahl der unabhängigen Komponenten des symmetrischen
        Tensors $g_{\mu\nu}$ in $D=4$ ($= D(D{+}1)/2$).
  \item $10 = 1+2+3+4$ (Dreieckszahl der Dimensionen).
  \item $10 = \binom{5}{2}$ (Paarzahl über den fünf Elementen
        $\{$vier Zyklen, ein Z$_3$-Fixpunkt$\}$ o.\,ä.).
\end{itemize}

\textbf{Warnung (R61-Disziplin):} Diese Lesarten sind post-hoc.
Der Auswahlraum kleiner ganzer Zahlen mit \enquote{schöner}
Kombinatorik ist groß -- dieselbe Lage wie bei den 32 glatten
2-3-5-Zahlen nahe $1/\xi$ (A-Serie): Ein Treffer ist strukturell
fast unvermeidlich und darum \emph{kein} Signal. Zulässig ist nur
eine Herleitung, die $n=10$ \emph{erzwingt}, bevor $H_0$ angeschaut
wird -- idealerweise mit einer zweiten, unabhängigen Konsequenz
desselben Mechanismus. Keine der drei Lesarten wird hier als
Ergebnis geführt. [SETZUNG-Kandidaten, nicht bookbar]

\section{B -- Der Ordnung-1-Faktor $\kappa$}

$\kappa = 2{,}12$ (Gl.~\eqref{eq:kappa}) muss aus der Geometrie
folgen (Kandidatenquelle: Krümmungs- bzw. Volumenfaktor der
kompakten T$^4$/Z$_3$-Struktur; in der Einstein-statischen Lösung
wäre der entsprechende Faktor exakt~1, in $\Lambda$CDM ist er
$3\,\Omega_\Lambda$). Bis zur Herleitung bleibt $\kappa$ offen und
wird \textbf{nicht} gefittet. Ein späterer \enquote{Treffer} durch
eine passende Konstantenkombination ohne Mechanismus wäre nach
R61-Maßstab zurückzuweisen.

\section{C -- Stabilität der statischen Konfiguration}

Störungsanalyse der statischen Konfiguration auf T$^4$/Z$_3$ bei
voll gültigen Feldgleichungen (Abschn.~\ref{sec:stabfrage}): Nachweis, dass die
instabilen Moden der klassischen Einstein-Lösung durch die kompakte
Topologie ausgeschlossen oder stabilisiert werden. Dies ist nach der
Übernahme der vollen Feldgleichungen (Abschn.~\ref{sec:links}) die
\emph{unausweichliche} dritte Herleitungspflicht -- ohne sie ist
die statische Lesart nur behauptet, nicht gesichert. [S]

\section{D -- Kompakte Zeitrichtung: KMS-Identifikation und
Kausalität}

Zweifache Pflicht aus dem Kapitel zur kompakten Zeitrichtung:
(i)~Vorwärtsbegründung, dass die kompakte Realzeitrichtung die
thermische KMS-Periodizität trägt (Gibbons--Hawking-Temperatur ohne
Horizont, Gl.~\eqref{eq:tgh}); (ii)~Nachweis, dass die
Feldkonfigurationen auf dem Zeitzyklus periodisch schließen
(Kausalitätsbedingung statt CTC-Paradoxie) -- in Kopplung an die
Modenzählung aus Anhang~\ref{app:stab}, da dieselben Randbedingungen
(periodisch/antiperiodisch) beide Fragen entscheiden. [S]

% ============================================================
\chapter{Die Zeit in den Feldgleichungen}
% ============================================================

\section{Innen, nicht außen [K]}

In $G_{\mu\nu} = (8\pi G/c^4)\,T_{\mu\nu}$ laufen die Indizes über
$0,1,2,3$ -- die $0$-Komponente \emph{ist} die Zeit ($x^0 = ct$).
Es gibt keine Gleichung \enquote{für den Raum} plus eine separate
Zeitentwicklung: Die Zeit ist eine der vier Richtungen der Metrik.
Konkret trägt $g_{00}$ den Zeittakt (Gravitationszeitdilatation),
die $g_{0i}$ die Mitführung (Lense--Thirring), die $g_{ij}$ den
Raum. Die Zwei-Hälften-Synthese (Dok.~307/308) ist exakt diese
Aufteilung: Zeitkopplung $\hat{=}$ $g_{00}$-Hälfte, fraktale
Wegverlängerung $\hat{=}$ $g_{rr}$-Hälfte.

Außerhalb der Gleichungen existiert dagegen \emph{keine} Zeit: kein
äußerer Uhrenparameter, gegen den sich das Ganze entwickelt. Die
Gleichungen sind kovariant -- jede Koordinatenwahl für
\enquote{Zeit} ist gleichberechtigt, keine physikalisch
ausgezeichnet. Zeit ist kein Behälter, sondern wird von der Lösung
erzeugt: Erst eine Metrik liefert Eigenzeiten entlang Weltlinien,
$d\tau^2 = -g_{\mu\nu}\,dx^\mu dx^\nu/c^2$.

\section{Zwangsbedingungen und das Problem der Zeit [K]}

Im ADM-3+1-Split zerfallen die zehn Gleichungen in sechs echte
Entwicklungsgleichungen und vier \textbf{Zwangsbedingungen},
darunter die Hamilton-Constraint $\mathcal{H} \approx 0$: Die
Gesamtenergie der Raumzeit ist null, es gibt keine äußere Uhr.
Kanonisch quantisiert wird daraus die Wheeler--DeWitt-Gleichung
$\hat{H}\Psi = 0$ -- eine Grundgleichung des Universums, in der
keine Zeitvariable mehr vorkommt. Das ist das bekannte
\enquote{Problem of Time}: Fundamental verschwindet die Zeit aus
der Beschreibung und muss relational -- aus Korrelationen zwischen
Feldern -- rekonstruiert werden.

\section{Anschluss an die FFGFT [S]}

Damit stehen drei Formulierungen derselben Richtung nebeneinander:

\begin{itemize}
  \item \textbf{ART:} Zeit entsteht aus der Metrik; kein äußerer
        Parameter.
  \item \textbf{Wheeler--DeWitt:} Fundamental gibt es keine
        Zeitvariable; Zeit ist relational.
  \item \textbf{FFGFT:} Zeit ist die Reziproke der Masse,
        $\tilde T \cdot m = 1$, verankert in der T$^4$-Struktur mit
        $\partial\text{T}^4 = \emptyset$ (Dok.~306) -- weshalb
        \enquote{was war bei $t=0$?} falsch gestellt ist.
\end{itemize}

Der Unterschied ist präzise benennbar: Die ART lässt offen,
\emph{woran} der Zeittakt hängt; die FFGFT gibt die Antwort -- an
der Masse. Das ist genau der Inhalt der $g_{00}$-Hälfte der
Zwei-Hälften-Synthese. Die Übernahme der vollen Feldgleichungen
(Abschn.~\ref{sec:links}) ist also auch auf der Zeitseite konsistent: Die FFGFT
widerspricht der relationalen Zeitstruktur der ART nicht, sondern
konkretisiert sie. Status [S]: Die Entsprechung
$\tilde T \cdot m = 1 \;\leftrightarrow\; g_{00}$-Sektor ist im
Schwachfeld gezeigt (Dok.~307); die Entsprechung zur vollen
Hamilton-Constraint ist eine Strukturaussage, kein Beweis.

% ============================================================
\chapter{Statik der Relation: Zeitdilatation als Massenlandkarte}
% ============================================================

\section{Die halbe Buchung [K]}

$\tilde T \cdot m = 1$ gilt grundsätzlich \emph{statisch}: Es ist
eine konstitutive Relation, keine Bewegungsgleichung. Die
Zeitdilatation wird dennoch fast immer dynamisch erzählt --
\enquote{die Uhr im Potential wird verlangsamt, das Photon verliert
beim Aufstieg Energie}. Diese Erzählung vergisst die Massenseite
der Buchung.

Die statische Gegenrechnung ist bekannt (Okun, Selivanov, Telegdi
2000): In einem statischen Gravitationsfeld ist die Photonfrequenz
auf dem gesamten Weg \emph{konstant} (Zeittranslationsinvarianz;
die Frequenz ist Erhaltungsgröße). Was sich unterscheidet, sind die
\emph{Atome} an verschiedenen Orten -- ihre Übergangsenergien, also
effektiv ihre Energie-/Massenniveaus im Potential. Die
Rotverschiebung entsteht beim \emph{Vergleich zweier verschiedener
Uhren}, nicht durch etwas, das dem Photon unterwegs
\enquote{passiert}.

In FFGFT-Sprache: $\tilde T(x)\cdot m(x) = 1$ gilt an jedem Ort;
die \enquote{Zeitdilatation} ist die statische \textbf{Landkarte}
$m(x)$, keine Dynamik. Wer dynamisch redet, ohne die Massenseite
mitzuführen, bucht nur die halbe Relation. Die $g_{00}$-Hälfte der
Zwei-Hälften-Synthese (Dok.~307) \emph{ist} diese Landkarte.

\section{Dieselbe Vergesslichkeit trägt die Expansionslesart [K]}

Auf kosmischer Skala existiert dazu ein publiziertes hartes
Resultat (Wetterich 2013, \emph{Universe without expansion}): Per
Feldredefinition (Weyl-Transformation) ist ein expandierendes
Universum mit konstanten Massen \textbf{exakt äquivalent} zu einem
statischen Raum mit anwachsenden Teilchenmassen. Beobachtbar sind
nur dimensionslose Verhältnisse; ob man $a(t)$ oder $m(t)$ laufen
lässt, ist Rahmenwahl, keine Physik. Die Expansionslesart wirkt nur
deshalb alternativlos, weil stillschweigend \enquote{Massen
konstant} gesetzt wird -- man vergisst, dass die Massen sich beim
Rahmenwechsel mitändern müssen.

Damit erhält die Entartung aus Dok.~267 ein externes Fundament:
Der Streit \enquote{expandiert oder nicht} ist in dieser Form nicht
entscheidbar. Entscheidbar ist nur, welcher Rahmen die
\emph{einfachere Setzungsstruktur} hat -- und das ist genau die
Frage dieses Dokuments ($\Lambda$ gefittet vs.\ Abschlussskala
abgeleitet).

\section{Konsequenz für die Lichtkurven-Streckung [S]}

Der historisch härteste Einwand gegen statische Modelle ist die
$(1+z)$-Streckung der Supernova-Lichtkurven: Klassisches
Tired Light scheitert daran, weil es nur die Photonseite ändert.
In der Massen-Lesart ist die Streckung dagegen automatisch:
$\tilde T \propto 1/m$ -- ändern sich die Massen, ändern sich
\emph{alle} Zeittakte reziprok mit; Wellenlänge und Uhrentakt sind
dieselbe Buchung. $\tilde T \cdot m = 1$ kann diesen Fehler
strukturell nicht machen. [S: qualitativ; der quantitative
Vorwärtsnachweis der vier Diskriminatoren -- Lichtkurven, $T(z)$,
Tolman, Alcock--Paczyński -- bleibt Programm.]

\section{Ehrliche Grenze}

Die Wetterich-Äquivalenz zeigt, dass die statische Lesart
\emph{möglich} und beobachtungsgleich ist -- nicht, dass sie
\emph{richtig} ist. Der FFGFT-Anspruch ist der stärkere: dass die
statische T$^4$-Fassung die \emph{natürliche} ist (keine
$\Lambda$-Setzung, Abschlussskala abgeleitet,
$\partial\text{T}^4 = \emptyset$), während die Expansionsfassung
dieselbe Physik in einen dynamischen Rahmen mit gefittetem
$\Lambda$ übersetzt. Dieser Anspruch steht und fällt mit den
offenen Punkten A--D -- er wird hier behauptet, nicht bewiesen.

% ============================================================
\chapter{Das Denken in Lichtwegen}
% ============================================================

\section{Der Bote ohne Eigenzeit [K]}

Fast unser gesamtes Bild von Zeit und Kosmos ist auf Lichtwegen
gebaut: Lichtkegel, Rotverschiebung, Distanzleitern,
Gleichzeitigkeit -- alles über Photonen definiert oder gemessen.
Dabei ist das Photon der eine Bote, der \emph{selbst keine Uhr
trägt}: Auf Nullgeodäten gilt $d\tau = 0$; entlang eines Lichtwegs
vergeht keine Eigenzeit; für $m = 0$ ist $\tilde T \cdot m = 1$
gar nicht definiert. Die Zeitstruktur der Welt über den
\emph{zeitlosen} Boten zu definieren, ist eine methodische
Umkehrung: Uhren sind Massen ($\nu = mc^2/h$, Compton-Uhr); Licht
ist bestenfalls das Lineal, nie die Uhr.

\section{Das Erbe des Operationalismus [K]}

Die Umkehrung hat einen klaren historischen Ursprung: Einstein 1905
definierte Gleichzeitigkeit \emph{operational über
Lichtsignale} -- eine brillante Messvorschrift, die aber
unbemerkt zur Ontologie wurde. Bemerkenswert ist, dass das
SI-System selbst die richtige Arbeitsteilung längst eingebaut hat:
Seit 1983 kommt die \emph{Länge} vom Licht ($c$ fixiert), die
\emph{Sekunde} aber von der Materie -- $\nu_\text{Cs}$, ein
Übergang in einem massiven Atom (vgl.\ Dok.~309: alle Einheiten
auf diesen einen Materie-Anker rückführbar). Das SI bucht Zeit auf
Materie und Länge auf Licht. $\tilde T \cdot m = 1$ macht diese
Arbeitsteilung explizit, statt sie im Lichtwege-Denken zu
verwischen.

\section{Folgen für die Kosmologie [S]}

Alle kosmologische Information erreicht uns auf \emph{einem}
Vergangenheits-Lichtkegel, und Licht transportiert nur
dimensionslose Verhältnisse. Daraus folgt zwingend: Wer
\emph{nur} in Lichtwegen denkt, kann die Rahmenfrage (Expansion
vs.\ Massenlauf, voriges Kapitel; Wetterich, Dok.~267) prinzipiell
nicht entscheiden -- die Entartung ist keine Schwäche der Daten,
sondern eine Folge des zeitlosen Boten. Entschieden wird, wo
\emph{Materie-Uhren} in der Ferne ablesbar sind: Die
$(1+z)$-Streckung der Supernova-Lichtkurven ist genau das -- eine
entfernte Materie-Uhr (die Physik des explodierenden Sterns), deren
Takt wir mit unserer vergleichen. $T(z)$-Messungen an
Molekülwolken ebenso: Thermometer aus Materie. Die
Diskriminatoren, die wirklich beißen, sind sämtlich
Materie-Ablesungen -- nicht Lichtwege.

\section{Einordnung in die Zwei-Hälften-Synthese [S]}

Die Synthese (Dok.~307) sagt dasselbe im Lokalen: Die Ablenkung
zerfällt in die Lichtweg-Hälfte (fraktale Wegverlängerung,
Dok.~182; $g_{rr}$) und die Materie-Uhr-Hälfte (Zeitkopplung
$\tilde T \cdot m = 1$; $g_{00}$). Einstein 1911 hatte nur die
Uhrenhälfte und bekam die halbe Ablenkung; wer heute nur in
Lichtwegen denkt, macht spiegelbildlich denselben Fehler auf der
anderen Seite. Vollständig ist erst die Buchung beider Hälften --
und die Kosmologie steht in dieser Hinsicht noch dort, wo die
Ablenkungsrechnung 1911 stand: eine Hälfte systematisch
unterbelichtet.

\section{Die Grenze bei $c$: relational, nicht innerlich [K/S]}

Nähert sich Materie der Lichtgeschwindigkeit, ändern sich Takt und
Masse gemeinsam: Für den ruhenden Beobachter geht die bewegte Uhr
um $\gamma$ langsamer, die Energie wächst um $\gamma$ --
$\tilde T \propto 1/\gamma$, $E \propto \gamma$, das Produkt bleibt
invariant. Das ist die kinetische Fassung von
$\tilde T \cdot m = 1$ und historisch exakt de~Broglies
\enquote{Harmonie der Phasen} (1924): innere Uhr um $\gamma$
verlangsamt, Phasenfrequenz um $\gamma$ erhöht, die Relation bleibt
geschlossen.

Daraus folgt zunächst das Gegenteil einer inneren Grenze: Weil sich
\emph{alle} Massen, Takte und Bindungsenergien eines Organismus
gemeinsam und reziprok ändern, merkt er intern nichts -- im
mitbewegten System ist die Physik exakt die gewohnte. Das
Relativitätsprinzip ist in dieser Lesart keine Setzung, sondern
Konsequenz der Ko-Änderung. Es gibt keine Geschwindigkeit, bei der
die Biologie \enquote{nicht mehr mitmacht}.

Die natürlichen Grenzen sitzen woanders, und alle drei sind hart
(Zahlen: \texttt{ffgft\_312\_grenze\_c.py}):

\begin{enumerate}
  \item \textbf{Beschleunigung, nicht Geschwindigkeit.} Der
        \emph{Übergang} kostet: dauerhaft verträgt ein Mensch
        $\sim$1--2\,g. Bei konstant 1\,g wächst $\gamma$ mit
        $\cosh(g\tau/c)$ -- etwa ein Jahr Eigenzeit pro Faktor $e$;
        die Grenze betrifft die Änderung, nie den Zustand.
  \item \textbf{Energie.} $\gamma$ divergiert: Ein 70-kg-Körper
        auf $\gamma = 10$ kostet $\sim 5{,}7\times10^{19}$~J --
        Größenordnung Weltjahresenergieverbrauch. Kein Verbot, eine
        praktische Mauer.
  \item \textbf{Die Umgebung -- die eigentliche Naturgrenze.} Die
        Ko-Änderung kehrt von außen zurück: Für den Reisenden
        blueshiftet die Umgebung um $\gamma$. Das interstellare
        Medium ($\sim$1~H-Atom/cm$^3$) wird zum Teilchenstrahl --
        bei $v = 0{,}99c$ trifft jedes Proton mit mehreren GeV, und
        die ungeschirmte Dosisleistung überschreitet die letale
        Dosis in Sekundenbruchteilen; bei sehr großem $\gamma$ wird
        selbst der CMB zur Röntgenquelle. Die Grenze lautet nicht
        \enquote{der Organismus verträgt die Geschwindigkeit
        nicht}, sondern: \emph{Er verträgt nicht, was die
        Geschwindigkeit aus seiner Umgebung macht.} Relational bis
        zuletzt -- im Geist von $\tilde T \cdot m = 1$: kein
        absoluter Zustand \enquote{schnell}, nur Beziehungen, und
        die Beziehung zum Medium wird tödlich, lange bevor
        irgendetwas Inneres versagt.

        Dieselbe Wand trifft \emph{unbelebte} Struktur, nur
        verschoben: Auch Elektronik ist nicht immun. Schon bei
        $0{,}2c$ ($\gamma \approx 1{,}02$) trifft jedes ISM-Proton
        mit $\sim$19~MeV -- mitten im Bereich maximaler
        Verlagerungsschäden in Halbleitern (Gitteratome werden von
        ihren Plätzen geschlagen, Transistoren degradieren
        kumulativ); dazu ionisierende Gesamtdosis und
        Single-Event-Effekte. Ein Nanogramm interstellaren Staubs
        trägt bei $0{,}2c$ die kinetische Energie einer
        Gewehrkugel ($\sim$1{,}8~kJ). Die Starshot-Studien (Hoang,
        Loeb et al.\ 2017) kommen genau dahin: Kante voraus
        fliegen, mm-dicke Opferschicht, die über die Flugzeit
        kontrolliert weggefressen wird -- und selbst dann ist das
        Überleben der Elektronik eine der offenen Kernfragen des
        Konzepts, keine gelöste. Der Unterschied zwischen Biologie
        und Elektronik ist nur die Dosisverträglichkeit
        ($\sim$5~Gy letal vs.\ $\sim$$10^2$--$10^5$~Gy
        funktionszerstörend je nach Härtung) und der
        Schadensmechanismus -- nicht das Prinzip.
\end{enumerate}

Die Grenze ist damit allgemeiner als eine Organismus-Grenze: Sie
trifft \textbf{jede strukturtragende Materie} -- alles, was
Information in geordneter Materie hält (Zellen, Kristallgitter,
Speicherzustände), ist der blueshifteten Umgebung ausgesetzt. In
der Sprache von A271: Strukturtragende Materie hält Zustände
(Ebene~2), und die blueshiftete Umgebung ist ein Löschangriff auf
diese Zustände -- die $c$-Grenze ist auch eine Grenze der
\emph{Informationserhaltung in bewegter Materie}. Wirklich immun
ist nur der Grenzfall selbst: das masselose, strukturlose Photon
mit $d\tau = 0$.

Damit schließt sich der Bogen zum Boten ohne Eigenzeit: Das Photon
($m = 0$, $d\tau = 0$) ist der Grenzfall, den Materie \emph{nie}
erreicht -- nicht weil eine innere Schranke bricht, sondern weil
die Energie divergiert und die Umgebung vorher zur Wand wird. $c$
ist für Massen keine erreichbare Geschwindigkeit, sondern die
Asymptote der Relation selbst ($\tilde T \to 0$ hieße
$m \to \infty$).

Zwei Klarstellungen gegen vorhersehbare Einwände von beiden Seiten:
Die Grenze \emph{bepreist} den relativistischen Reisetraum
($\gamma \gg 1$ als Vehikel der Zeitdilatation), sie
\emph{verbietet} langsame interstellare Fahrt nicht -- bei
$0{,}1$--$0{,}2c$ sind Dosisleistungen um Größenordnungen milder
und Abschirmung denkbar. Und sie gilt unabhängig davon, wem sie
gefällt: Maßstab ist die deklarierte Annahme
($n = 1$/cm$^3$, ungeschirmt, 5~Gy), nicht die Erwünschtheit des
Ergebnisses -- derselbe Standard, der Wunschresultate \emph{für}
kühne Projekte zurückweist, schützt sie auch vor
Unmöglichkeits-Pathos ohne Buchung.

% ============================================================
\chapter{Die kompakte Zeitrichtung}
% ============================================================

\section{Die Konsequenz [K|P20]}

Wenn die Zeit eine der vier Metrikrichtungen ist (voriges Kapitel)
und die FFGFT auf T$^4$ mit $\partial\text{T}^4 = \emptyset$ lebt,
dann ist auch die \textbf{Zeitrichtung kompakt}: Die Feldgleichungen
gelten nicht auf $\mathbb{R}\times\text{T}^3$, sondern auf dem
vollen T$^4$. Bei isotropem Torus (alle vier Zyklen gleich lang,
$L^{*} = 2\pi R_H$) hat der Zeitzyklus automatisch die Länge

\begin{equation}
  \tau_\text{Zyklus} \;=\; \frac{L^{*}}{c}
  \;=\; \frac{2\pi}{H_0}
  \;=\; 2{,}90\times10^{18}~\text{s}
  \;\approx\; 91{,}9~\text{Gyr}.
  \label{eq:tauzyk}
\end{equation}

Keine neue Setzung: dieselbe Leiterstufe $\xi^{10}$, vierte Rolle
derselben Skala (nach $\Lambda^{*}$, $T_w$, $\hbar H_0$). Das
Energiequant des Zeitzyklus, $2\pi\hbar/\tau_\text{Zyklus}
= \hbar H_0$, ist identisch mit dem Impulsquant der Raumzyklen am
C1-Minimum -- Isotropie der vier Richtungen auch auf Quantenebene.

\section{Kompakte Zeit und Temperatur [S]}

In der Feldtheorie ist Periodizität in der \emph{euklidischen} Zeit
mit Periode $\tau$ äquivalent zu einer Temperatur
$T = \hbar/(k_B\tau)$ (KMS-Bedingung, Matsubara-Formalismus). Trägt
man die Zeitzyklus-Periode ein -- mit dem KMS-üblichen Faktor
$2\pi$, d.\,h. $\tau_\text{eukl} = \tau_\text{Zyklus}$ als
$2\pi/H_0$ --, folgt

\begin{equation}
  \boxed{\;T \;=\; \frac{\hbar H_0}{2\pi k_B}
  \;=\; 2{,}63\times10^{-30}~\text{K}\;}
  \label{eq:tgh}
\end{equation}

Das ist \textbf{exakt die Gibbons--Hawking-Temperatur} des
de-Sitter-Horizonts -- hier jedoch ohne Horizont und ohne Expansion:
Die Temperatur folgt aus der Kompaktheit der Zeitrichtung selbst.

Status ehrlich [S]: Die KMS-Äquivalenz gilt für periodische
\emph{Imaginärzeit} (nach Wick-Rotation). Die Identifikation der
kompakten Realzeitrichtung der FFGFT mit dieser thermischen
Periodizität ist eine Strukturaussage, kein Beweis -- sie gehört zu
den Vorwärtspflichten (Punkt~D).

\section{Vereinheitlichung mit der $a_0$-Kette [S]}

Setzt man die Unruh-Temperatur $T_U = \hbar a/(2\pi c\,k_B)$ gleich
der Temperatur~\eqref{eq:tgh}, folgt

\begin{equation}
  a \;=\; c\,H_0 \;=\; 6{,}49\times10^{-10}~\text{m/s}^2,
\end{equation}

also genau der Anker der $a_0$-Herleitung aus Dok.~308
(Übergangsskala des Trägheitsregimes, $cH_0/2\pi
= 1{,}03\times10^{-10}$~m/s$^2$ gegen
$a_0^\text{RAR} \approx 1{,}2\times10^{-10}$). Was in Dok.~308 zwei
Argumente waren (Unruh-Verankerung; $\xi^{10}$-Kette), wird hier
eines: \textbf{Die Kompaktheit der Zeitrichtung ist die strukturelle
Quelle der Unruh-Verankerung.}

\section{Der Preis: geschlossene Zeitkurven [offen]}

Kompakte Zeitrichtung heißt formal: Weltlinien kehren nach
$\tau_\text{Zyklus}$ zu ihrem Ausgangswert der Zeitkoordinate
zurück -- in der ART das klassische Problem geschlossener Zeitkurven
(CTC). Zwei Entschärfungen sind formulierbar, beide [S]:

\begin{enumerate}
  \item \textbf{Empirisch:} Die Periode $\sim 92$~Gyr übersteigt
        jede beobachtbare Zeitspanne; ein Beobachtungskonflikt
        existiert nicht.
  \item \textbf{Relational:} In der Lesart des vorigen Kapitels ist
        die Zeitkoordinate kein Behälter, in den man
        \enquote{zurückkehrt}. Periodische Koordinatenzeit ist nur
        dann paradox, wenn die Feldkonfigurationen sich \emph{nicht}
        periodisch schließen. Die Kausalitätsbedingung lautet also:
        \textbf{Periodizität der Feldkonfigurationen auf dem
        Zeitzyklus} -- auf T$^4$ genau die natürliche Randbedingung.
        Dieselbe Wahl (periodisch vs.\ antiperiodisch pro Feld)
        bestimmt zugleich die Modenzählung der Kippgrenze in
        Anhang~\ref{app:stab} (C1$'$): Kausalität und Stabilität hängen an
        \emph{derselben} Randbedingungsstruktur.
\end{enumerate}

Ein Beweis ist keines von beiden; die CTC-Frage geht als Punkt~D in
den Pflichtenkatalog.

% ============================================================
\chapter{Relativbewegung und das Torus-Ruhesystem}
% ============================================================

Wir bewegen uns: die Erde um die Sonne, die Sonne um die Galaxis,
die Lokale Gruppe Richtung Großer Attraktor. Wie passt das in eine
statische kompakte Lesart? Die Antwort hat eine lokale und eine
globale Hälfte, und beide sind gerechnet
(\texttt{ffgft\_312\_relativbewegung.py}).

\section{Lokal: Bewegung auf der Landkarte [K]}

Alle Bahnbewegungen sind nichtrelativistisch:

\begin{center}
\begin{tabular}{lrr}
\toprule
Bewegung & $v$ [km/s] & $\gamma-1$ \\
\midrule
Erdrotation (Äquator) & $0{,}5$ & $1{,}2\times10^{-12}$ \\
Erde um Sonne & $29{,}8$ & $4{,}9\times10^{-9}$ \\
Sonne um Galaxis & $233$ & $3{,}0\times10^{-7}$ \\
Sonne rel.\ CMB & $369{,}8$ & $7{,}6\times10^{-7}$ \\
Lokale Gruppe rel.\ CMB & $620$ & $2{,}1\times10^{-6}$ \\
\bottomrule
\end{tabular}
\end{center}

Der Uhrentakt eines bewegten Körpers ist das Produkt zweier
Buchungen: statische Landkarte $m(x)$ (Kapitel \enquote{Statik der
Relation}) mal kinetischer Faktor $1/\gamma$ (Ko-Änderung). Genau so
rechnet GPS: Potentialterm plus Geschwindigkeitsterm. Bahnbewegung
ist Bewegung \emph{auf} der Landkarte, kein Einwand gegen sie. Die
Hierarchie schließt sich zudem vektoriell: Sonne-rel.-CMB minus
Galaxisbahn ergibt $552$~km/s für die Milchstraße relativ zum CMB
-- Literaturwert $\sim550$.

\section{Global: die Topologie zeichnet ein System aus [K]}

Auf kompakter Topologie gilt Lorentz-Invarianz \emph{lokal}
uneingeschränkt, aber Boosts sind keine globalen Symmetrien mehr:
Das Zwillingsparadoxon hat auf dem Torus eine absolute Antwort --
der im Torus-Ruhesystem verbleibende Zwilling altert am meisten,
unterscheidbar über die \emph{Windungszahl} der Weltlinie.
Quantitativ: Bei $\beta = 0{,}99$ dauert eine Umrundung
$92{,}9$~Gyr Koordinatenzeit bei $13{,}1$~Gyr Eigenzeit
(Defizit $79{,}8$~Gyr). Die Auszeichnung ist strukturell absolut,
praktisch aber inert -- jede Umrundung dauert mindestens
$\tau_\text{Zyklus} = 92$~Gyr. Kein Äther: Nichts Materielles
reibt; die Auszeichnung ist topologisch, keine lokale Messung kann
sie detektieren, nur globale Marker.

\section{Der CMB-Dipol als Kandidat [S]}

Ein solches ausgezeichnetes System wird beobachtet: Der CMB-Dipol
liefert mit $\Delta T/T = 1{,}2336\times10^{-3}$ exakt
$v = 369{,}8$~km/s -- die Doppler-Ablesung eines Ruhesystems, auf
$0{,}03\,\%$ konsistent mit der direkten Messung. In der
Expansionslesart ist das \enquote{nur} das mitbewegte System, eine
Konvention. In der kompakten Lesart ist es der Kandidat für das
\textbf{Torus-Ruhesystem} selbst: Alle Bewegungen der Hierarchie
sind Geschwindigkeiten relativ zur Struktur, der Dipol ihre Summe.
Auch die kompakte Zeitrichtung hängt daran: $\tau_\text{Zyklus}$,
das $\hbar H_0$-Quant und die KMS-Temperatur sind im Torus-System
definiert; ein bewegter Beobachter liest die Modenstruktur
dopplerverschoben ab -- wieder der Dipol.

Status ehrlich: Die Identifikation Torus-System $=$ CMB-System ist
naheliegend und konsistent, aber eine Identifikation [S], keine
Ableitung.

\textbf{Topologische Periode und Lichtumlauf sind zu trennen.}
Die Periode des Zeitzyklus ist eine \emph{topologische} Größe,
$\tau_c = L^{*}/c = 91{,}9$~Gyr. Ein Photon bräuchte für eine
Umrundung wegen der fraktalen Wegverlängerung (Dok.~313)
$1{,}35\,\%$ länger ($93{,}2$~Gyr). Für KMS, Unruh und
Gibbons--Hawking zählt die \emph{Feldperiodizität}, nicht die
Photonlaufzeit: $T = \hbar H_0/(2\pi k_B)$, das Quant
$\hbar H_0$ und $a = cH_0$ bleiben daher unverändert.

\section{Prüfsignaturen und Selbstkonsistenz [S]}

\begin{enumerate}
  \item \textbf{Matched-Circle-Versatz.} Bei kompakter Topologie
        nahe der Beobachtungsgrenze ($L^{*}/D_\text{LSS} \approx
        1{,}01$) müssten Kreissignaturen um den Dipolvektor um
        $\sim\beta$~rad $= 4{,}2$~Bogenminuten versetzt sein --
        eine spezifische, prinzipiell zugängliche
        Richtungssignatur -- \emph{sofern} die Streuwand den
        Halbzyklus überschreitet. Mit der fraktalen Wegkorrektur
        (Dok.~313) rückt sie auf $0{,}979\,\pi$ und schneidet
        sich nicht mehr selbst; dann existieren keine Matched
        Circles und diese Signatur entfällt. Sie gilt also nur in
        der glatt gerechneten Fassung; als Prüffläche verbleibt
        dann die Nicht-Monotonie am Antipoden (Dok.~313).
  \item \textbf{Drift.} Unsere Bewegung durch die Struktur beträgt
        $34{,}8$~Mpc pro Zeitzyklus ($0{,}12\,\%$ von $L^{*}$) und
        $5{,}2$~Mpc über die Rekombinations-Lichtlaufzeit
        ($0{,}019\,\%$): Die Eigenbewegung erzeugt keinen
        Selbstwiderspruch der statischen Lesart.
\end{enumerate}

Kurzfassung: Relativbewegung bleibt lokal relativ (SRT
unangetastet) und wird global relational zur Struktur. Was in der
Expansionslesart Konvention ist (das mitbewegte System), ist hier
eine Strukturaussage mit benannter Prüfsignatur.

% ============================================================
\chapter{Warum Einstein $\tilde T \cdot m = 1$ nicht sah}
% ============================================================

\emph{Historisch-methodische Einordnung; keine physikalische
Buchung. Die Frage ist ernst gemeint: Die Zutaten der Relation
stammen von Einstein selbst.}

\section{Die Zutaten lagen auf seinem Tisch}

$E = mc^2$ (1905) und $E = h\nu$ (1905, Lichtquanten) ergeben
zusammen

\begin{equation}
  \nu \;=\; \frac{mc^2}{h}
  \qquad\Longleftrightarrow\qquad
  T \cdot m \;=\; \frac{h}{c^2} \;\;(\to 1
  \text{ in natürlichen Einheiten}).
\end{equation}

Jede Masse \emph{ist} eine Frequenz, also eine Uhr -- das ist
$\tilde T \cdot m = 1$ in Rohform, aus zwei Einstein-Formeln.
Ausgesprochen hat es \textbf{de~Broglie 1924}: Die
\enquote{innere Uhr} des Teilchens mit $\nu = mc^2/h$ war der Kern
seiner Dissertation. Einstein hat die Arbeit begutachtet und
gelobt (\enquote{er hat einen Zipfel des großen Schleiers
gelüftet}) -- aber alle, auch de~Broglie selbst, machten daraus die
\emph{Wellenlänge} und ließen die \emph{Uhr} fallen. Die
Compton-Zeit $\hbar/(mc^2)$ galt fortan als abgeleitete Skala,
nicht als das, was Zeit \emph{ist}.

\section{Vier Gründe}

\textbf{1. Sein Programm lief andersherum.} Für Einstein war die
Geometrie das Primäre und die Masse eine \emph{Quelle} -- rechts in
der Gleichung, als $T_{\mu\nu}$. Die Denkfigur \enquote{Masse
krümmt die Bühne} verhindert die Denkfigur \enquote{Masse
\emph{ist} der Takt}. $\tilde T \cdot m = 1$ macht die Masse zum
Konstitutiven der Zeit selbst -- die Umkehrung, die sein
Feldprogramm nicht vorsah. Die Ironie: 1911 hatte er \emph{nur} die
Zeitkopplung (die halbe Ablenkung -- die $g_{00}$-Hälfte der
Zwei-Hälften-Synthese). Er stand auf der $T$-Seite, bevor er zur
Raumgeometrie überging.

\textbf{2. In der Relation steckt $\hbar$.} Ausgeschrieben lautet
sie $T = \hbar/(mc^2)$ -- sie verheiratet Quantenkonstante und
Relativität auf fundamentalster Ebene. Einstein hielt die
Quantentheorie für vorläufig und wollte sie \emph{ersetzen}, nicht
mit ihr fusionieren. Eine Ontologie, deren Grundrelation $\hbar$
enthält, war für ihn kein Fundament, sondern ein Symptom des
Unfertigen. Die FFGFT-Ausgangshaltung -- beide Theorien
wohlbestätigt, beide unvollständig, kein fundamentaler Widerspruch
möglich, also dualisieren -- ist exakt die Haltung, die er nicht
einnehmen konnte.

\textbf{3. Die Einheiten-Klarheit fehlte.} $\tilde T \cdot m = 1$
setzt voraus, dass $\hbar$ und $c$ als bloße Umrechnungsfaktoren
durchschaut sind -- dass Zeit und Masse keine unabhängigen
Dimensionen tragen. Diese Sicht wurde erst mit der Praxis
natürlicher Einheiten in der Teilchenphysik selbstverständlich und
wirklich konsequent erst 2019, als das SI alle Einheiten auf
Konstanten-Fixierung umstellte (vgl.\ Dok.~309: alle SI-Einheiten
auf $\nu_\text{Cs}$ rückführbar). Zu Einsteins Zeit waren Sekunde
und Gramm ontologisch verschiedene Dinge; eine Reziprozität
zwischen ihnen sah aus wie ein Kategorienfehler.

\textbf{4. Es sieht nicht aus wie ein Gesetz.} Einstein suchte
Feldgleichungen -- Dynamik. $\tilde T \cdot m = 1$ ist keine
Bewegungsgleichung, sondern eine konstitutive Relation: in der
Terminologie von \emph{Ordnung ohne Rangordnung} ein
\textbf{Grundsatz}, kein Gesetz. Solche Aussagen wirken trivial
(die Compton-Zeit kannte jeder), solange man sie nicht ontologisch
nimmt. Der Schritt ist nicht die Formel, sondern der
\emph{Statuswechsel} der Formel.

\section{Späte Bestätigung und Schluss}

Die Uhr-Lesart wurde erst 2013 operativ bestätigt: Das
Compton-Uhr-Experiment (Lan, Müller et al., Berkeley) maß eine
Masse \emph{als Frequenznormal} -- Masse und Zeittakt operativ
dasselbe. Da war die Relation 88~Jahre alt.

Der ehrliche Schluss: keine Blindheit, sondern Programmtreue.
Einstein sah die Hälfte, die in sein Bild passte (Masse krümmt den
Zeittakt), und nicht die Reziproke (der Zeittakt \emph{ist} die
Masse), weil sie $\hbar$ zum Fundament gemacht hätte -- den einen
Stein, den er nicht verbauen wollte. Die FFGFT verbaut genau diesen
Stein -- und übernimmt im Gegenzug seine Gleichungen unverändert
(Abschn.~\ref{sec:links}): Arbeitsteilung statt Widerlegung.

% ============================================================
\chapter{Statusübersicht}
% ============================================================

\begin{center}
\begin{tabular}{lp{6.5cm}l}
\toprule
Aussage & Inhalt & Status \\
\midrule
Identifikation & $\Lambda^{*} = 1/R_H^2$, kein neuer Parameter
  & [K|P20] \\
Dimensionslose Form & $\Lambda^{*}\bar\lambda_e^2 = (\pi/2)^2\xi^{20}$
  & [K|P20] \\
Strukturzerlegung & $10^{-122}$ als $\xi^{20}\cdot(l_P/\bar\lambda_e)^2$
  & [K|P20] \\
Reduktion & Feinabstimmung $\to$ P20 (eine offene Stelle statt zwei)
  & [S] \\
Lokaler Sektor & $\Lambda^{*}$ lokal automatisch vernachlässigbar
  & [K] \\
Feldgleichungen & voll übernommen (Lovelock); FFGFT liefert
  $\Lambda$-, $G$-Wert, Lösungswahl & [S] \\
Zeit & innen (Metrik, $g_{00}$); relational; FFGFT konkretisiert:
  $\tilde T \cdot m = 1$ & [K]/[S] \\
Statik & Zeitdilatation als Massenlandkarte $m(x)$ (Okun et al.);
  Expansion vs.\ Massenlauf rahmenäquivalent (Wetterich) & [K] \\
Lichtwege & Photon ohne Eigenzeit ($d\tau=0$): Licht = Lineal,
  Masse = Uhr; Diskriminatoren sind Materie-Ablesungen & [K]/[S] \\
Grenze bei $c$ & relational: innerlich keine ($\gamma$-Ko-Änderung);
  hart: Beschleunigung, Energie, Umgebung & [K]/[S] \\
Zeitzyklus & $\tau = 2\pi/H_0 = 91{,}9$~Gyr; Quant $= \hbar H_0$
  & [K|P20] \\
Temperatur & $\hbar H_0/(2\pi k_B)$ = Gibbons--Hawking, ohne
  Horizont; Quelle der Unruh-/$a_0$-Verankerung & [S] \\
Kausalität & CTC $\to$ Periodizität der Felder auf dem Zeitzyklus
  & offen (D) \\
Ruhesystem & Topologie zeichnet global ein System aus; Kandidat:
  CMB-Dipol ($369{,}8$~km/s exakt); Prüfsignatur $4{,}2'$-Versatz
  nur ohne fraktale Wegkorrektur (Dok.~313) & [K]/[S] \\
Stabilität & offen; Kandidat: kompakte Topologie als Randbedingung
  & offen (C) \\
$\kappa = 2{,}12$ & Ordnung-1-Faktor & offen \\
Exponent 10 & Vorwärtsherleitung & offen (P20) \\
\bottomrule
\end{tabular}
\end{center}

\medskip
\noindent\textbf{Registrierung:} Die Aufnahme der Abschlussskala und
der hier benannten offenen Punkte in Dok.~190 (Korrekturregister,
append-only) sowie etwaige Folgeanpassungen in Dok.~308/309 erfolgen
separat und sind nicht Teil dieses Dokuments.


% ============================================================
\appendix
\chapter{Anhang: Berechnungen zur Stabilitätsfrage (C1--C3)}
\label{app:stab}
% ============================================================

\section*{Worum es geht}
% ============================================================

Abschnitt~\ref{sec:stabfrage} (Punkt~C) lässt die Stabilitätsfrage offen: Bei
voll gültigen Feldgleichungen ist der Torusradius dynamisch -- die
Eddington-Frage kehrt als Modulus-Frage wieder (kompaktifiziert oder
ausgerollt?), strukturell dieselbe Falle wie der
Decompactification-Runaway im T$^6$/Z$_3$-Review (Quílez Zamora).
Dieses Beiblatt rechnet drei Dinge durch:

\begin{itemize}
  \item \textbf{C1} -- Radius-Stabilisierung als
        Windungs-/Impulsmoden-Gleichgewicht (Spielzeugmodell,
        Brandenberger--Vafa-Typ), \emph{invers} gerechnet: Welche
        Windungsspannung erzwingt das Minimum bei $L^{*} = 2\pi R_H$?
  \item \textbf{C2} -- Eddington-Kontrast: Zeitkonstante der
        klassischen Instabilität vs.\ Vorzeichen von $V''$ im
        Torus-Modell.
  \item \textbf{C3} -- Ordnungsprüfung gegen CMB-Topologiegrenzen.
\end{itemize}

Alle Zahlen: \texttt{ffgft\_312C\_stabilisierung.py}. Status
durchgehend [K|P20]\,$\times$\,[S] (Spielzeugmodell: ein Zyklus,
feste Quantenzahlen). Keine Fits -- $T_w$ wird invers bestimmt und
anschließend auf der $\xi$-Leiter \emph{gebucht}; der Treffer auf
Stufe $\xi^{20}$ ist Ergebnis, nicht Eingabe.

% ============================================================
\section{C1 -- Das Windungs-/Impulsmoden-Gleichgewicht}
% ============================================================

\subsection{Modell}

Effektives Potential für die Zyklenlänge $L$ eines T$^4$-Zyklus:

\begin{equation}
  V(L) \;=\; \underbrace{N_w\,T_w\,L}_{\text{Windung}}
  \;+\; \underbrace{N_k\,\frac{2\pi\hbar c}{L}}_{\text{Impulsquanten}}
  \;\;(+\,\text{Casimir, klein}).
  \label{eq:pot}
\end{equation}

Windungsmoden tragen Energie $\propto L$ (Spannung mal Länge;
topologisch, nicht stetig deformierbar), Impulsquanten
$\propto 1/L$. Die Konkurrenz erzeugt ein Minimum bei

\begin{equation}
  L_{\min} = \sqrt{\frac{2\pi N_k \hbar c}{N_w T_w}},
  \qquad V''(L_{\min}) > 0.
\end{equation}

Dies ist der Mechanismus, der dem T$^6$/Z$_3$-Modell im Review
fehlte: Dort gab es keinen $\propto L$-Term, der das Ausrollen
aufhält.

\subsection{Inverse Rechnung [K|P20]}

Forderung $L_{\min} = L^{*} = 2\pi R_H = 4\bar\lambda_e/\xi^{10}$
(mit $N_w = N_k = 1$) bestimmt die Windungsspannung:

\begin{equation}
  \boxed{\;T_w \;=\; \frac{2\pi\hbar c}{L^{*2}}
  \;=\; \frac{\hbar c\,\Lambda^{*}}{2\pi}
  \;=\; \frac{\pi}{8}\,\frac{E_e}{\bar\lambda_e}\,\xi^{20}
  \;=\; 2{,}63\times10^{-79}~\text{N}\;}
  \label{eq:tw}
\end{equation}

Beide Identitäten sind exakt (numerisch auf $10^{-16}$ bestätigt).
Zwei Befunde:

\begin{enumerate}
  \item \textbf{Eine Skala, zwei Rollen.} Die Spannung, die das
        Minimum bei $R_H$ erzwingt, ist direkt proportional zur
        Abschlussskala: $T_w = \hbar c\,\Lambda^{*}/(2\pi)$.
        Abschlussskala (Dok.~312) und Stabilisierungsspannung sitzen
        auf \emph{derselben} Leiterstufe $\xi^{20}$ -- es wird keine
        zweite unabhängige Skala benötigt. [S]
  \item \textbf{Selbstkonsistenz am Minimum.} Das Impulsquant am
        Minimum ist exakt
        \[
          \frac{2\pi\hbar c}{L^{*}} = \frac{\hbar c}{R_H}
          = \hbar H_0 = 2{,}28\times10^{-52}~\text{J}
          = 1{,}4\times10^{-33}~\text{eV},
        \]
        und am Minimum gilt Gleichverteilung
        $E_\text{Windung} = E_\text{Impuls}$ (Brandenberger--Vafa-%
        Eigenschaft). Die Energieskala des stabilisierten Zyklus
        \emph{ist} die Hubble-Energie -- nichts Neues muss gesetzt
        werden.
\end{enumerate}

\subsection{Casimir-Korrektur und Kippgrenze}

Ein masseloser periodischer Boson-Freiheitsgrad trägt
$E_C = -\pi\hbar c/(6L)$ bei ($|E_C|/E_\text{Impuls} = 1/12$ pro
Freiheitsgrad). Verschiebung des Minimums:

\begin{center}
\begin{tabular}{lll}
\toprule
$N_\text{Casimir}$ (netto) & $L_{\min}/L^{*}$ & Befund \\
\midrule
0  & $1{,}000$ & unverändert \\
4  & $0{,}817$ & Minimum bleibt, verschoben \\
24 & --        & kein Minimum: Kollaps \\
\bottomrule
\end{tabular}
\end{center}

\textbf{Ehrliche Grenze:} Ab $N_\text{Casimir} \ge 12\,N_k$ kippt
das Modell. Die Vorwärtspflicht schließt also die Modenzählung des
realen T$^4$/Z$_3$-Spektrums ein (Bosonen minus Fermionen,
periodisch/antiperiodisch) -- nicht nur die Spannungsherleitung.

% ============================================================
\section{C2 -- Eddington-Kontrast [K]}
% ============================================================

Die klassische Einstein-Lösung (1917) hat kein Minimum, sondern
einen Balancepunkt; Störungen wachsen exponentiell mit

\begin{equation}
  \tau \;=\; \frac{1}{c\sqrt{\Lambda^{*}}} \;=\; \frac{R_H}{c}
  \;=\; 4{,}6\times10^{17}~\text{s} \;=\; 14{,}6~\text{Gyr}.
\end{equation}

Die Instabilität ist also keine akademische Sorge: Sie kippt die
Konfiguration auf der Hubble-Zeitskala. Kein statisches Modell kann
sich hinter einer langsamen Instabilität verstecken.

Im Torus-Modell gilt dagegen $V''(L^{*}) > 0$: Störungen des Radius
oszillieren, statt wegzulaufen. Die Oszillationsfrequenz hängt an
der Modulus-Trägheit $M_\text{mod}$ (offen); entscheidend für die
Eddington-Frage ist hier allein das \emph{Vorzeichen} von $V''$.

% ============================================================
\section{C3 -- Ordnungsprüfung: CMB-Topologiegrenze [S]}
% ============================================================

\begin{equation}
  L^{*} = 28{,}2~\text{Gpc}
  \qquad\text{vs.}\qquad
  D_\text{LSS} \approx 27{,}8~\text{Gpc}
  \qquad\Rightarrow\qquad
  \frac{L^{*}}{D_\text{LSS}} = 1{,}01\text{--}1{,}02.
\end{equation}

\textbf{Mit fraktaler Wegkorrektur.} $D_\text{LSS}$ ist ein
\emph{Lichtweg}, $L^{*}$ eine topologische Größe -- nach der
Zuordnungsregel (Dok.~313) greift $K_\text{frak} = 1-100\xi$
hier. Damit verschiebt sich das Verhältnis auf
$L^{*}/D_\text{LSS} = 1{,}03$: Die Streuwand liegt
\emph{deutlicher} unter dem Halbzyklus. Das ist mit den
Nullresultaten der Topologiesuchen verträglicher, kostet aber die
Matched-Circle-Signatur (Kapitel Relativbewegung).

Matched-Circles-Suchen (Planck) schließen Zyklenlängen deutlich
\emph{unter} dem Durchmesser der letzten Streufläche aus; $L^{*}$
liegt etwa 1--2\,\% \emph{darüber}: knapp konsistent, hart an der
Beobachtungsgrenze. Das ist eine Stärke, keine Schwäche -- die
kompakte Lesart ist damit \textbf{falsifizierbar}, nicht beliebig.

\textbf{Vorbehalt:} Distanzmaße sind in der statischen Lesart neu zu
interpretieren (Dok.~267-Entartung); dies ist eine Ordnungsprüfung,
kein Endtest.

% ============================================================
\section{Ergebnis und Vorwärtspflichten}
% ============================================================

\begin{center}
\begin{tabular}{lp{7.2cm}l}
\toprule
& Befund & Status \\
\midrule
C1 & $T_w = \hbar c\Lambda^{*}/(2\pi) = (\pi/8)(E_e/\bar\lambda_e)\xi^{20}$;
     Impulsquant am Minimum $= \hbar H_0$ & [K|P20]$\times$[S] \\
C1' & Kippgrenze: Minimum verschwindet ab
     $N_\text{Casimir} \ge 12\,N_k$ & [K] (Modell) \\
C2 & Einstein-statisch: $\tau = R_H/c = 14{,}6$~Gyr;
     Torus: $V'' > 0$ & [K] \\
C3 & $L^{*}/D_\text{LSS} = 1{,}01$--$1{,}02$ (fraktal $1{,}03$) -- an der Grenze,
     falsifizierbar & [S] \\
\bottomrule
\end{tabular}
\end{center}

\medskip
\noindent\textbf{Was C1 leistet und was nicht:} Die inverse Rechnung
zeigt, dass die Stabilisierung bei $R_H$ \emph{keine neue Skala}
verlangt -- die Abschlussskala trägt beide Rollen, und die
Energieskala am Minimum ist $\hbar H_0$. Sie zeigt \emph{nicht},
dass die T$^4$/Z$_3$-Windungsdynamik diese Spannung tatsächlich
liefert. Die Vorwärtspflichten sind damit präzise:

\begin{enumerate}
  \item $T_w$ aus der Windungsdynamik auf T$^4$/Z$_3$ ableiten
        (Zielwert Gl.~\eqref{eq:tw} -- vor der Rechnung versiegelt,
        R61-Disziplin).
  \item Modenzählung des realen Spektrums (Kippgrenze
        $N_\text{Casimir} < 12\,N_k$ prüfen).
  \item Modulus-Trägheit $M_\text{mod}$ für die
        Oszillationsfrequenz (C2 quantitativ).
\end{enumerate}

\noindent\textbf{Registrierung:} Aufnahme in Dok.~190 erfolgt
separat.


\end{document}
